% ============================================================
% PREPIUM RENAISSANCE MOCK — IMPORT TEMPLATE
% ============================================================
% HOW TO USE:
%   1. Copy EVERYTHING below into the "Full LaTeX" box in
%      Admin > Renaissance > Builder > Bulk Import.
%   2. Click "Preview Split" to check the questions.
%   3. Click "Import" — answers from the Answer Key at the
%      bottom are filled in automatically.
%
% RULES THE IMPORTER UNDERSTANDS:
%   - Questions: use \item inside the numbered enumerate below.
%   - Options:    use \begin{enumerate}[label=\Alph*.] ... \item ...
%   - Math:       use $...$ (inline) or $$...$$ (display).
%   - Diagrams:   use \begin{circuitikz}...\end{circuitikz} or
%                 \begin{tikzpicture}...\end{tikzpicture} (inside
%                 \begin{center}...\end{center}).
%   - Answer key: a heading line "Answer Key" / "Answers" at the
%                 END of the paper, then one line per answer,
%                 e.g.  1 & B   or   1. B   or   1  B.
%   - Comments (lines starting with %) are ignored.
%
% NOTE: You can write "A. ..." instead of the option enumerate too.
% ============================================================

\begin{enumerate}[label=\textbf{\arabic*.}]

% ---------- Example 1: plain MCQ ----------
\item What is the current in a wire when a charge of $6\,\mathrm{C}$ flows in $3\,\mathrm{s}$?

\begin{enumerate}[label=\Alph*.]
    \item $0.5\,\mathrm{A}$
    \item $2.0\,\mathrm{A}$
    \item $3.0\,\mathrm{A}$
    \item $18.0\,\mathrm{A}$
\end{enumerate}

% ---------- Example 2: MCQ with a circuit diagram ----------
\item Consider the circuit below.

\begin{center}
\begin{circuitikz}
\draw
(0,0) to[battery1,l=$12\,V$] (0,3)
      to[R,l=$6\,\Omega$] (4,3)
      -- (4,0)
      -- (0,0);
\end{circuitikz}
\end{center}

What is the current flowing through the resistor?

\begin{enumerate}[label=\Alph*.]
    \item $0.5\,\mathrm{A}$
    \item $2.0\,\mathrm{A}$
    \item $6.0\,\mathrm{A}$
    \item $72\,\mathrm{A}$
\end{enumerate}

% ---------- Example 3: MCQ with a graph ----------
\item The graph shows potential difference $V$ against current $I$ for an ohmic conductor.

\begin{center}
\begin{tikzpicture}[scale=0.9]
\draw[->] (0,0) -- (5,0) node[right] {$I$};
\draw[->] (0,0) -- (0,4) node[above] {$V$};
\draw[thick] (0,0) -- (4,3.2);
\end{tikzpicture}
\end{center}

What does the gradient of the graph represent?

\begin{enumerate}[label=\Alph*.]
    \item Current
    \item Resistance
    \item Charge
    \item Power
\end{enumerate}

% ---------- Example 4: theory / open question ----------
\item State Ohm's law and give the equation $V = IR$.

% ---------- Example 5: MCQ using letters instead of enumerate ----------
\item Which equation represents Ohm's law?

A. $V = IR$
B. $V = I/R$
C. $R = VI$
D. $I = VR$

\end{enumerate}

% ============================================================
% ANSWER KEY  (keep at the very end of the paper)
% Question number  &  Correct letter, one per line.
% Works as a table OR plain lines:  1. B   /   1 & B   /   1  B
% ============================================================
\newpage

\section*{Answer Key}

\begin{tabular}{|c|c|}
\hline
\textbf{Question} & \textbf{Answer} \\
\hline
1 & B \\
\hline
2 & B \\
\hline
3 & B \\
\hline
4 & --- \\
\hline
5 & A \\
\hline
\end{tabular}
